% GENERATED FILE. Edit canonical lessons, then run npm run generate:books.
% canonical-source-hash: fnv1a64:10fd6d50a10e46f7
% canonical-lessons: GE-C395-passwort, GE-C395-einloggen, GE-C395-zugang, GE-C395-datenschutz, GE-C395-verschluesseln

\chapter{Sicher im Netz}
\label{ch:ge-sicher-im-netz}
\hlchaptermodality{395}

\begin{chapteropening}
\textbf{By the end of this chapter:} \emph{I can talk about passwords, logging in, access and data protection.}

Complete the last lesson of chapter 395: I can talk about passwords, logging in, access and data protection.
\end{chapteropening}

\section[das Passwort]{das Passwort — the password}

\label{lesson:GE-C395-passwort}



\begin{quote}
Before the new one: say the German for to post, then the German for to comment on.
\end{quote}

\subsection*{You'll want to know: das Passwort}
\textbf{das Passwort} — \textquotedblleft{}the password\textquotedblright{}.

\begin{quote}
Ich habe mein Passwort schon wieder vergessen.
\end{quote}

\begin{quote}
\emph{I have forgotten my password yet again.}
\end{quote}

Plural: \emph{die Passwörter}. Also \emph{das Kennwort}.

\begin{cousinweb}
\emph{Passwort} was borrowed from English \emph{password} and fitted to German \emph{Wort} \textquotedblleft{}word\textquotedblright{}, inherited from Proto-Germanic \emph{wurdą} (EtymDB).
\end{cousinweb}

\begin{grammarlens}[title={What you've built}]
One more word for telling, explaining and describing at B1.
\end{grammarlens}

\subsection*{Your turn}
\begin{itemize}
\raggedright
  \item \emph{Say it:} \emph{das Passwort}
  \item \emph{Say it:} \emph{das Passwort}, once more
  \item \emph{Say it:} a sentence of your own with \emph{das Passwort}
  \item \emph{Recall:} say the German for to post, then the German for to comment on, then say \emph{das Passwort} again
\end{itemize}

\begin{tcolorbox}[breakable,colback=teal!4,colframe=teal!35!black,title={Before you move on}]
What does das Passwort mean? (\textquotedblleft{}The password\textquotedblright{}.) Say it once more.
\end{tcolorbox}
\section[einloggen]{einloggen — to log in (reflexive, sich einloggen)}

\label{lesson:GE-C395-einloggen}



\begin{quote}
Before the new one: say the German for to comment on, then the German for the password.
\end{quote}

\subsection*{You'll want to know: einloggen}
\textbf{einloggen} — \textquotedblleft{}to log in (reflexive, sich einloggen)\textquotedblright{}.

\begin{quote}
Ohne Passwort kann ich mich nicht einloggen.
\end{quote}

\begin{quote}
\emph{I cannot log in without a password.}
\end{quote}

Reflexive and separable: \emph{ich logge mich ein}, \emph{ich habe mich eingeloggt}. Also \emph{sich anmelden}.

\begin{cousinweb}
\emph{einloggen} takes English \emph{log in} and gives it the German prefix \emph{ein-} \textquotedblleft{}in\textquotedblright{}. English \emph{log in} is usually traced to the ship's log, the daily book named after the wooden log once used to measure a ship's speed; EtymDB has no entry for the German verb.
\end{cousinweb}

\begin{grammarlens}[title={What you've built}]
One more word for telling, explaining and describing at B1.
\end{grammarlens}

\subsection*{Your turn}
\begin{itemize}
\raggedright
  \item \emph{Say it:} \emph{einloggen}
  \item \emph{Say it:} \emph{einloggen}, once more
  \item \emph{Say it:} a sentence of your own with \emph{einloggen}
  \item \emph{Recall:} say the German for to comment on, then the German for the password, then say \emph{einloggen} again
\end{itemize}

\begin{tcolorbox}[breakable,colback=teal!4,colframe=teal!35!black,title={Before you move on}]
What does einloggen mean? (\textquotedblleft{}To log in (reflexive, sich einloggen)\textquotedblright{}.) Say it once more.
\end{tcolorbox}
\section[der Zugang]{der Zugang — the access; the way in}

\label{lesson:GE-C395-zugang}



\begin{quote}
Before the new one: say the German for the password, then the German for to log in.
\end{quote}

\subsection*{You'll want to know: der Zugang}
\textbf{der Zugang} — \textquotedblleft{}the access; the way in\textquotedblright{}.

\begin{quote}
Im Hotel gibt es kostenlosen Zugang zum Internet.
\end{quote}

\begin{quote}
\emph{The hotel has free internet access.}
\end{quote}

Plural: \emph{die Zugänge}. \emph{Zugang zu} + dative.

\begin{cousinweb}
\emph{Zugang} joins \emph{zu} \textquotedblleft{}to\textquotedblright{} to \emph{Gang} \textquotedblleft{}going, path\textquotedblright{}, inherited from Old High German \emph{gang} and Proto-Germanic \emph{gangaz} (EtymDB): the way to something.
\end{cousinweb}

\begin{grammarlens}[title={What you've built}]
One more word for telling, explaining and describing at B1.
\end{grammarlens}

\subsection*{Your turn}
\begin{itemize}
\raggedright
  \item \emph{Say it:} \emph{der Zugang}
  \item \emph{Say it:} \emph{der Zugang}, once more
  \item \emph{Say it:} a sentence of your own with \emph{der Zugang}
  \item \emph{Recall:} say the German for the password, then the German for to log in, then say \emph{der Zugang} again
\end{itemize}

\begin{tcolorbox}[breakable,colback=teal!4,colframe=teal!35!black,title={Before you move on}]
What does der Zugang mean? (\textquotedblleft{}The access; the way in\textquotedblright{}.) Say it once more.
\end{tcolorbox}
\section[der Datenschutz]{der Datenschutz — data protection, privacy}

\label{lesson:GE-C395-datenschutz}



\begin{quote}
Before the new one: say the German for to log in, then the German for the access.
\end{quote}

\subsection*{You'll want to know: der Datenschutz}
\textbf{der Datenschutz} — \textquotedblleft{}data protection, privacy\textquotedblright{}.

\begin{quote}
Wegen des Datenschutzes darf ich Ihnen die Nummer nicht geben.
\end{quote}

\begin{quote}
\emph{Because of data protection I am not allowed to give you the number.}
\end{quote}

No plural. \emph{Die Daten} — the data.

\begin{cousinweb}
\emph{Datenschutz} joins \emph{Daten} \textquotedblleft{}data\textquotedblright{}, from Latin \emph{datum} \textquotedblleft{}something given\textquotedblright{} (EtymDB), to \emph{Schutz} \textquotedblleft{}protection\textquotedblright{}, Middle High German \emph{schuz} (EtymDB).
\end{cousinweb}

\begin{grammarlens}[title={What you've built}]
One more word for telling, explaining and describing at B1.
\end{grammarlens}

\subsection*{Your turn}
\begin{itemize}
\raggedright
  \item \emph{Say it:} \emph{der Datenschutz}
  \item \emph{Say it:} \emph{der Datenschutz}, once more
  \item \emph{Say it:} a sentence of your own with \emph{der Datenschutz}
  \item \emph{Recall:} say the German for to log in, then the German for the access, then say \emph{der Datenschutz} again
\end{itemize}

\begin{tcolorbox}[breakable,colback=teal!4,colframe=teal!35!black,title={Before you move on}]
What does der Datenschutz mean? (\textquotedblleft{}Data protection, privacy\textquotedblright{}.) Say it once more.
\end{tcolorbox}
\section[verschlüsseln]{verschlüsseln — to encrypt, to encode}

\label{lesson:GE-C395-verschluesseln}



\begin{quote}
Before the new one: say the German for the access, then the German for data protection.
\end{quote}

\subsection*{You'll want to know: verschlüsseln}
\textbf{verschlüsseln} — \textquotedblleft{}to encrypt, to encode\textquotedblright{}.

\begin{quote}
Die Nachrichten werden automatisch verschlüsselt.
\end{quote}

\begin{quote}
\emph{The messages are encrypted automatically.}
\end{quote}

Inseparable: \emph{es verschlüsselt, es hat verschlüsselt}. Built on \emph{der Schlüssel} — the key.

\begin{cousinweb}
\emph{verschlüsseln} is built with \emph{ver-} on \emph{Schlüssel} \textquotedblleft{}key\textquotedblright{}, inherited from Old High German \emph{sluzzil} and Proto-Germanic \emph{slutilaz} (EtymDB), of the family of \emph{schließen} \textquotedblleft{}to lock\textquotedblright{}.
\end{cousinweb}

\begin{grammarlens}[title={What you've built}]
One more word for telling, explaining and describing at B1.
\end{grammarlens}

\subsection*{Your turn}
\begin{itemize}
\raggedright
  \item \emph{Say it:} \emph{verschlüsseln}
  \item \emph{Say it:} \emph{verschlüsseln}, once more
  \item \emph{Say it:} a sentence of your own with \emph{verschlüsseln}
  \item \emph{Recall:} say the German for the access, then the German for data protection, then say \emph{verschlüsseln} again
\end{itemize}

\begin{tcolorbox}[breakable,colback=teal!4,colframe=teal!35!black,title={Before you move on}]
What does verschlüsseln mean? (\textquotedblleft{}To encrypt, to encode\textquotedblright{}.) Say it once more.
\end{tcolorbox}
